\begin{lemma}
The incident core preserves Tashkinov order:
$t(G^{\circ})=t(G)$.
\end{lemma}
\begin{proof}
Put $k=\chi'(G)-1$.  By
\noderef{IncidentCoreChromaticIndex}, this is also
$\chi'(G^{\circ})-1$, so both maxima in \noderef{TashkinovOrder} use the
same palette size.

Let $n$ be attainable in the defining maximum for $G$.  Thus there are an
edge identity $e_0$, a partial coloring $\varphi$, and an ordered
tree-sequence $T$ such that $(G,e_0,\varphi)$ is a $k$-triple, $T$ is a
Tashkinov tree, and its vertex set has cardinality $n$.
By \noderef{IncidentCoreKTriple}, the same edge identity and coloring form a
$k$-triple on $G^{\circ}$.  The forward half of
\noderef{IncidentCoreTreeTransport} supplies a certified core tree with the
same indexed edge identities and vertex-set cardinality $n$.  Hence every
attainable value for $G$ is attainable for $G^{\circ}$.

Conversely, begin with any value attainable for $G^{\circ}$.  The reverse
halves of \noderef{IncidentCoreKTriple} and
\noderef{IncidentCoreTreeTransport} give a $k$-triple and a certified tree
of $G$ having the same vertex-set cardinality.  Thus the two sets of
attainable natural numbers coincide.  The finite range bounds in
\noderef{TashkinovOrder} remove no attainable value, because a tree's
injective vertex list has cardinality at most its ambient finite vertex
type; the two transport directions also verify the corresponding bound in
the other graph.  Both definitions adjoin $0$ and take the maximum of the
same attainable set, so their maxima, and hence their Tashkinov orders, are
equal.
\end{proof}
